\section*{Chapter \ref{ch:iv:algorithms-and-data-structures} --- Algorithms and Data Structures}

\begin{solution}{iq:iv:algorithms-and-data-structures:1}
One pass with a hash map from size to index: for each size $x$, look up $t - x$ among the sizes already seen; if present,
return the pair, else record $x$. $O(n)$ time and memory. The brute force is $O(n^2)$; sorting with two pointers is $O(n\log n)$
and loses the original indices unless they are carried.
\iqlookfor{the hash-map complement and the comparison with the alternatives.}
\end{solution}

\begin{solution}{iq:iv:algorithms-and-data-structures:2}
A hash set of identifiers seen; return the first one already present. $O(n)$ time, $O(n)$ memory. If the stream does not fit:
partition identifiers by a hash into files that do, then check each partition (two passes); or use a Bloom filter as a first
filter and verify its positives exactly, accepting a small false-positive rate on the first pass.
\iqlookfor{the set solution and an external-memory plan.}
\end{solution}

\begin{solution}{iq:iv:algorithms-and-data-structures:3}
Two pointers, taking the earlier head each time: $O(n + m)$. For ties, decide and state a rule that makes the merge stable and
deterministic, such as tape A before tape B, or a secondary key (sequence number, venue); the chapter's code keeps A first and
is checked against sorting on 300 random cases.
\iqlookfor{the linear merge and an explicit, deterministic tie rule.}
\end{solution}

\begin{solution}{iq:iv:algorithms-and-data-structures:4}
Track the lowest price so far and the best difference: $O(n)$, $O(1)$ memory. The answer is 0 if prices only fall (no trade).
Example: $7, 1, 5, 3, 6, 4$ gives 5.
\iqlookfor{the running-minimum scan and the no-trade case.}
\end{solution}

\begin{solution}{iq:iv:algorithms-and-data-structures:5}
Binary search for the first element not less than $t$: \texttt{bisect\_left} in Python, \texttt{std::lower\_bound} in C++;
$O(\log n)$. It returns the length of the array when every timestamp is earlier, which the caller must handle.
\iqlookfor{lower bound, its return value at the end, and the library call.}
\end{solution}

\begin{solution}{iq:iv:algorithms-and-data-structures:6}
A \omterm{def:iv:algorithms-and-data-structures:families}{monotonic deque} of indices with decreasing prices (\cref{lst:iv:algorithms-and-data-structures:slidingmax}): pop smaller
prices from the back before pushing, pop the front when it leaves the window, and read the maximum at the front. Each index
enters and leaves once: $O(n)$ overall, $O(k)$ memory. Tested against the quadratic oracle.
\iqlookfor{the deque invariant and the amortised argument.}
\end{solution}

\begin{solution}{iq:iv:algorithms-and-data-structures:7}
Two heaps: a max-heap for the lower half and a min-heap for the upper half, balanced to within one element; the median is at the
top. Removing the value that leaves the window is done lazily: record it and discard it when it reaches a top. $O(\log k)$ per
trade. With integer sizes below 1\,000, keep a count per size and walk to the middle, or keep a Fenwick tree over sizes to find
the median in $O(\log 1000)$: bounded values beat comparisons.
\iqlookfor{the two-heap median with deletion, and exploiting a bounded domain.}
\end{solution}

\begin{solution}{iq:iv:algorithms-and-data-structures:8}
Sum volume per symbol in a hash map, then keep a min-heap of size $k$ over the totals: $O(n + s\log k)$ for $s$ symbols. With an
unbounded stream and bounded memory, use a heavy-hitters sketch (Space-Saving or a Count-Min sketch with a heap), which
guarantees the frequent symbols with bounded error.
\iqlookfor{counting plus a bounded heap, and a streaming sketch when memory binds.}
\end{solution}

\begin{solution}{iq:iv:algorithms-and-data-structures:9}
Sweep over the sorted endpoints, $+1$ at a start and $-1$ at an end, with ends before starts at equal times if orders are open on
$[s, e)$; the largest running total is the answer (\cref{lst:iv:algorithms-and-data-structures:sweep}). $O(n\log n)$. For
$(0,5), (1,3), (3,6), (5,7)$ the answer is 2.
\iqlookfor{the event sweep and the tie convention that encodes the interval's openness.}
\end{solution}

\begin{solution}{iq:iv:algorithms-and-data-structures:10}
Sort by start; extend the last merged interval while the next one starts before or at its end, otherwise open a new one.
$O(n\log n)$. Say whether touching intervals merge (the chapter's code merges them).
\iqlookfor{sort and merge, with the touching-interval convention stated.}
\end{solution}

\begin{solution}{iq:iv:algorithms-and-data-structures:11}
A cycle multiplies money when the product of its rates exceeds one, that is when the sum of $-\log(\text{rate})$ along it is
negative. Run Bellman--Ford from a virtual source (all distances 0) for $n - 1$ rounds and check whether any edge can still be
relaxed: $O(n^3)$ for $n$ currencies. In practice add a tolerance for fees and floating point, and recover the cycle from the
predecessor pointers to trade it. The chapter's example turns 1 into about 1.064.
\iqlookfor{the log transformation, negative-cycle detection and the practical caveats.}
\end{solution}

\begin{solution}{iq:iv:algorithms-and-data-structures:12}
Dynamic programming over (day, trades used, holding or not): $\text{hold}_j = \max(\text{hold}_j, \text{free}_{j-1} - p)$ and
$\text{free}_j = \max(\text{free}_j, \text{hold}_j + p - \text{fee})$, $O(nk)$ time and $O(k)$ memory. On $3, 2, 6, 5, 0, 3$
with $k = 2$ and no fee the answer is 7; the chapter's code checks the recursion against exhaustive search on small inputs.
\iqlookfor{the state definition and the recursion.}
\end{solution}

\begin{solution}{iq:iv:algorithms-and-data-structures:13}
An ordered map from price to quantity (\cref{lst:iv:algorithms-and-data-structures:levels}): add and cancel in $O(\log n)$,
best price in $O(1)$ at the map's end. With integer ticks in a known narrow band, an array indexed by tick offset gives $O(1)$
add and cancel, and a bitmap of non-empty levels finds the best price with a few word scans; the array must be recentred when
prices drift. Real books also keep orders within a level in time priority (a linked list per level) and an index from order
identifier to order (One Quant Book~13, chapter~19).
\iqlookfor{the tree solution, the array-of-ticks alternative, and what a production book adds.}
\end{solution}

\begin{solution}{iq:iv:algorithms-and-data-structures:14}
Welford's update adds a value: $n \mathrel{+}= 1$, $\delta = x - \bar x$, $\bar x \mathrel{+}= \delta/n$, $M_2 \mathrel{+}=
\delta(x - \bar x)$; removing a value reverses it: $n \mathrel{-}= 1$, $\delta = y - \bar x$, $\bar x \mathrel{-}= \delta/n$,
$M_2 \mathrel{-}= \delta(y - \bar x)$; the variance is $M_2/(n - 1)$. $O(1)$ per update and far more stable than keeping $\sum x$
and $\sum x^2$, which cancel catastrophically when the mean is large. Over very long streams, recompute the window exactly now and
then to stop rounding errors accumulating.
\iqlookfor{Welford with removal, and the reason naive sums fail.}
\end{solution}
